\documentclass[10pt,a4paper]{amsart}
\usepackage[margin=19mm]{geometry}
\usepackage{amsmath,amssymb,mathtools}
\usepackage[scr=rsfs,cal=euler]{mathalfa}
\usepackage{xfrac}
\usepackage[unicode,hidelinks,bookmarks=false]{hyperref}
\newcommand{\ie}{i.e.\ }
\newcommand{\cf}{cf.\ }
\newcommand{\resp}{respectively\ }
\newcommand{\Subsection}{\subsection}
\providecommand{\nobreakdash}{\nobreak\hbox{-}}
\setlength{\emergencystretch}{2em}
\multlinegap0pt
\usepackage{algorithm}
\usepackage{algpseudocode}
\usepackage{amsbsy}
\usepackage{bm}
\usepackage{setspace}
\usepackage{float}
\newcommand{\sys}{\ensuremath{p}}
\newcommand{\dist}{\mathcal{F}}
\newcommand{\designdist}{\ensuremath{G}}
\newcommand{\elem}[2]{\ensuremath{#1[#2]}}
\newcommand{\fea}{\ensuremath{\tau}}
\newcommand{\spa}[1]{\text{{$\mathcal{#1}$}}}
\newcommand{\knn}{\ensuremath{k}}
\newcommand{\des}{\ensuremath{\bm{x}}}
\newcommand{\Ran}{\ensuremath{\bm{Y}}}
\newcommand{\rss}{\ensuremath{r}}
\newcommand{\rssCap}{\ensuremath{R}}
\newcommand{\rssVec}{\ensuremath{\mathcal{R}}}
\newcommand{\ball}[2]{\ensuremath{B(#1,#2)}}
\newcommand{\dis}{\lambda}
\newcommand{\diff}{\eta}
\newcommand{\goodnum}{N}
\newcommand{\goodnumLower}{n}
\newcommand{\vect}[1]{\ensuremath{\bm{#1}}}
\newcommand{\ratio}{\ensuremath{o}}
\newcommand{\altAlpha}{\ensuremath{\alpha_A}}
\newcommand{\altDelta}{\ensuremath{\delta_A}}
\newcommand{\desCap}{\ensuremath{\bm{X}}}
\newcommand{\design}{\ensuremath{\mathcal{D}}}
\newcommand{\test}{\ensuremath{\mathcal{U}}}
\newcommand{\testP}{\ensuremath{\bm{Z}}}
\newcommand{\classHat}{\widehat{\jmath}}
\newcommand{\dataU}{\mathrm{O}}
\newcommand{\dataL}{\mathrm{o}}
\newcommand{\dataEstU}{\mathrm{S}}
\newcommand{\dataEstL}{\mathrm{s}}
\newcommand{\dataEstSamp}{\mathrm{V}}
\newcommand{\dataS}{\mathcal{O}}
\newcommand{\deltaDiff}{\Delta}
\newcommand{\deltaDiffQ}[2]{\Delta(#1)_{1-#2}}
\newcommand{\deltaDiffQHat}[2]{\widehat{\Delta}_{1-#2}}
\newcommand{\deltaDiffQHatEst}[2]{\widehat{\overline{\Delta}}_{1-#2}}
\newcommand{\quant}{\ensuremath{\Delta(\vect{X})}}
\newcommand{\classHatU}{\widehat{J}}
\newcommand{\metric}{d}
\newcommand{\argmin}{\text{argmin}}
\newcommand{\argmax}{\text{argmax}}
\newcommand{\orderToRegular}{j^+}
\newcommand{\classHatLU}{\widehat{\overline{J}}}
\newcommand{\classHatL}{\widehat{\overline{\jmath}}}
\newcommand{\deltaDiffVecEst}{\widehat{\overline{\Delta}}^\dagger}
\newcommand{\deltaDiffVecExp}{\overline{\Delta}^\dagger}
\newcommand{\closestL}{\overline{d}}
\newcommand{\rssVecL}{\rho}
\newcommand{\deltaDiffQTwo}{\widehat{\Delta}}
\newcommand\givenbase[1][]{\:#1\lvert\:}
\newcommand\sgiven{\givenbase[\delimsize]}
\newcommand\Average{E\Basics}
\newtheorem{lemma}{Lemma}
\begin{document}
\title[A Proof of the Convergence a Lipschitz Upper Confidence Boun]{A Proof of the Convergence a Lipschitz Upper Confidence Bound}
\author{Gregory Keslin}
\date{}
\maketitle
\noindent\emph{Source and licence.} A Proof of the Convergence a Lipschitz Upper Confidence Bound. Version arXiv:2608.03081v1 (CC BY 4.0). This material is distributed under the Creative Commons Attribution 4.0 International licence, http://creativecommons.org/licenses/by/4.0/, as recorded in the arXiv OAI licence metadata for this exact version and shown on the abstract page https://arxiv.org/abs/2608.03081v1. Full rights notice: licenses/arxiv-2608.03081-CC-BY-4.0.txt.\par\smallskip
\noindent\textit{Editorial scope.} Counted unit: the original \texttt{lemma} environment \texttt{lemma: ulln} at source lines 140--144 together with its complete proof at lines 145--211; the delivered document keeps source lines 110--212 so that every referenced label and every citation resolves inside it. Native editable source is the paper's own concatenated TeX; the AMS wrapper is editorial formatting only. No publisher class, upstream script or unrelated artwork is redistributed.
\section{Lipschitz Estimation Proofs}
In the prior section, the Plausible Optima problem was discussed when the Lipschitz constant of $\mu$ is known. However, it is rare that the Lipschitz constant is known. Therefore, in such settings, it must be estimated. 

If not all covariates have been sampled, it is not possible to estimate an upper bound of the Lipschitz constant since the Lipschitz constant is a supremum across the entire covariate space. However, it is feasible to construct a $1-\alpha$ lower confidence bound of the Lipschitz constant. This document proves the consistency of this lower confidence bound. Because the lower confidence bound is consistent, it may make sense to use it as a point estimator of the Lipschitz constant. By increasing $\alpha$, the lower confidence bound can be made to overestimate the Lipschitz constant with higher probability. Therefore, in some sense, for large $\alpha$, we can upper bound the Lipschitz constant. 

Define $n(k,\vect{x})$ to be the function which maps the number of replications generated from $\vect{x}_i$ for design size $k$. We assume $n(k,\vect{x})$ is larger than 3 for all $k$ and $\vect{x}$. Fix $u\in \mathbb{R}^+$. Let $\overline{M}_\lambda=\bigcup_{0\leq \lambda^\prime \leq \lambda }M_\lambda$. Define $\overline{M}^u_\lambda=\{\phi \in \overline{M}_\lambda|\sup_{\vect{x}\in \spa{X}}|\phi(\vect{x})|\leq u\}$. Set $d^1(\phi,\hat{\mu},\hat{\sigma}^2)=\sum_{i=1}^k\sqrt{n(k,\des_i)}|\hat{\mu}_i-\phi(\des_i)|/\hat{\sigma}_i$.

We are interested in a $1-\alpha$ lower confidence bound of the Lipschitz constant of $\mu$, $\lambda^\star$. We define this confidence bound as the inversion of a hypothesis test,
$$\hat{\lambda}_\alpha=\inf_{\lambda} \text{ s.t } \inf_{\phi \in \overline{M}_\lambda} d^1(\phi,\hat{\mu},\hat{\sigma}^2) \leq D_\alpha.$$
By definition,
$$\mathbb{P}\left (\lambda^\star>\hat{\lambda}_\alpha \right )\geq 1-\alpha.$$
We wish to prove that as $k \to \infty$, for $\lambda<\lambda^\star$, the hypothesis test above has power $1$,
$$\lim_{k \to \infty} \mathbb{P}\left (\inf_{\phi \in \overline{M}_\lambda} d^1(\phi,\hat{\mu},\hat{\sigma}^2)>D_\alpha \right )=1.$$

To do this, we prove a series of lemmas. The first is a Uniform Law of Large Numbers for the smaller class of Lipschitz functions, $\overline{M}_\lambda^u$, when the number of replications generated at each covariate are bounded. We then extend the result to the class we are interested in, $\overline{M}_\lambda$, when the number of replications generated at each covariate are bounded. Finally, we drop the requirement that the number of replications generated at each covariate are bounded.

The following is a list of assumptions that we assume on the problem:

1. Distribution: Let $Y_j(\vect{x})$ be a replication generated from covariate $\vect{x}$. Then, $Y_j(\vect{x})\sim N(\mu(\vect{x}),\sigma^2(\vect{x}))$. 

2. Bounded Variances: $0<\underline{\sigma}\leq \sigma(\vect{x})\leq \overline{\sigma}$.

3. Uniform Monotone Lipschitz Continuity of $n(k,x)$: There exists $\lambda^\prime$ such that for all $k$, $\max_{\vect{x},\vect{x}^\prime}n(k,\vect{x})-n(k,\vect{x}^\prime)\leq \lambda^\prime s(\vect{x},\vect{x}^\prime)$ and for each $\vect{x}$, $n(\cdot,\vect{x})$ is a non-decreasing function.

4. Minimum Replication Count: For all $k$, $\vect{x} \in \spa{X}, n(k,\vect{x})>3$.

5. Compactness: $\spa{X}$ is compact.

6. Design Distribution: Each $\desCap_i$ is independently generated from probability measure $\pi$.


\begin{lemma}
\label{lemma: ulln}
 A Uniform Law of Large Numbers: Suppose $\sup_{k,\vect{x} \in \spa{X}} n(k,\vect{x}) \leq N$. For all $u$ and $\epsilon>0$,
 $$\lim_{k \to \infty}\mathbb{P}\left (\sup_{\phi \in \overline{M}^u_\lambda}\frac{|d^1(\phi,\hat{\mu},\hat{\sigma}^2)-\mathbb{E}\left ( d^1(\phi,\hat{\mu},\hat{\sigma}^2)\right )|}{k}>\epsilon\right )=0.$$
\end{lemma}

\begin{proof}
    By definition, $\overline{M}^u_\lambda$ is a set of equicontinuous uniformly bounded functions. Let $r(\cdot)$ be the metric induced by the sup norm of continuous functions, $r(\phi,\phi^\prime)=||\phi-\phi^\prime||_\infty=\sup_{x \in \spa{X}}|\phi(x)-\phi^\prime(x)|.$ Then, since $\spa{X}$ is compact, for the standard metric topology induced by $r(\cdot)$, by the Arzel\`a-Ascoli Theorem, $\overline{M}^u_\lambda$ is compact. 

    Notice that for all $\delta>0$, if $r(\phi,\phi^\prime)<\delta$, 
    \begin{equation*}
        \begin{split}
            \frac{\left |d^1(\phi,\hat{\mu},\hat{\sigma}^2)-d^1(\phi^\prime,\hat{\mu},\hat{\sigma}^2)\right |}{k}&=\left | \sum_{i=1}^k\frac{\sqrt{n(k,\desCap_i)}(|\hat{\mu}_i-\phi(\desCap_i)|-|\hat{\mu}_i-\phi^\prime(\desCap_i)|)}{k \hat{\sigma}_i} \right |\\
            & \leq \left | \sum_{i=1}^k\frac{\sqrt{n(k,\desCap_i)}(|\hat{\mu}_i-\phi(\desCap_i)-\hat{\mu}_i+\phi^\prime(\desCap_i)|)}{k \hat{\sigma}_i} \right |\\
            & \leq \delta \left | \sum_{i=1}^k\frac{\sqrt{n(k,\desCap_i)}}{k \hat{\sigma}_i} \right |\\
            & \leq \frac{N\delta}{\bar{\sigma}} \sum_{i=1}^k \frac{1}{k\chi_{i,n(k,\desCap_i)-1}},
        \end{split}
    \end{equation*}
    where $\chi_{i,n(k,\desCap_i)-1}^2$, when conditioned on $\desCap_i=\des_i$, is, for each $i$, an independent chi-squared random variable with $n(k,\des_i)-1$ degrees of freedom. Because $n(k,\des_i)>3$, $\mathbb{E}(1/\chi_{i,n(k,\desCap_i)-1}^2)=\mathbb{E}(1/(n(k,\desCap_i)-3))\leq 1$. Therefore, $\mathbb{E}(1/\chi_{i,n(k,\desCap_i)-1})\leq 2$. Thus, since $(1/\chi_{i,n(k,\desCap_i)-1})_{i=1}^k$ is an i.i.d sequence, by a triangular S.L.L.N., there exists $k_0$ such that,
    $$\mathbb{P}\left ( \sup_{k\geq k_0} \sum_{i=1}^k \frac{1}{k\chi_{i,n(k,\desCap_i)-1}} > 3\right )\leq \frac{\epsilon}{2}.$$
    Define the event $A=\{\sup_{k\geq k_0} \sum_{i=1}^k \frac{1}{k\chi_{i,n(k,\desCap_i)-1}} \leq 3\}$.
    
    (The remainder of the proof follows standard arguments for uniform strong laws of large numbers; see, e.g., Newey \cite{newey1991uniform}.)

    Since $\overline{M}^u_\lambda$ is compact, there exists a finite covering of $\overline{M}^u_\lambda$, $g_1,g_2,\ldots,g_t$ such that for all $\phi \in \overline{M}^u_\lambda$, there exists $g_j$ where $r(g_j,\phi)\leq \epsilon\underline{\sigma}/(9N)$. For each $g_j$, since $n(k,\des_i)\leq N$ almost surely and $||g_j||_\infty\leq u$, $\sup_k \mathbb{E}\left ((\sqrt{n(k,\des_i)}(\hat{\mu}_i-g_j(\des_i)))^2/\hat{\sigma}^2_i\right )<\infty$. Thus, by a triangular law of large numbers (see Durrett \cite{durrett2019probability}, for example), there exists $k_j$ such that for all $k>k_j$,
    $$\mathbb{P}\left (\frac{|d^1(g_j,\hat{\mu},\hat{\sigma}^2)-\mathbb{E}\left ( d^1(g_j,\hat{\mu},\hat{\sigma}^2)\right )|}{k}>\frac{\epsilon}{3}\right )<\frac{\epsilon}{2t}.$$
    
    Let $k^\star=\max_{j=1}^t k_j$. Consider $k\geq k^\star$. Fix arbitrary $\phi \in \overline{M}^u_\lambda$. Choose $g_j$ such that $r(g_j,\phi)\leq \epsilon/(9N)$. By the Triangle Inequality,
    \begin{equation*}
        \begin{split}
           \left |d^1(\phi,\hat{\mu},\hat{\sigma}^2)-\mathbb{E}\left ( d^1(\phi,\hat{\mu},\hat{\sigma}^2)\right )\right |&\leq \left |d^1(\phi,\hat{\mu},\hat{\sigma}^2)-d^1(g_j,\hat{\mu},\hat{\sigma}^2) \right | +\left | d^1(g_j,\hat{\mu},\hat{\sigma}^2)-\mathbb{E}\left ( d^1(g_j,\hat{\mu},\hat{\sigma}^2)\right ) \right |\\
           & +\left |\mathbb{E}\left ( d^1(g_j,\hat{\mu},\hat{\sigma}^2)\right )-\mathbb{E}\left ( d^1(\phi,\hat{\mu},\hat{\sigma}^2)\right ) \right |.
        \end{split}
    \end{equation*}
    From Jensen's inequality,
    \begin{equation*}
        \begin{split}
            \frac{1}{k}\left |\mathbb{E}\left ( d^1(g_j,\hat{\mu},\hat{\sigma}^2)\right )-\mathbb{E}\left ( d^1(\phi,\hat{\mu},\hat{\sigma}^2)\right ) \right |&\leq  \frac{1}{k}\mathbb{E}\left (\left | d^1(g_j,\hat{\mu},\hat{\sigma}^2)- d^1(\phi,\hat{\mu},\hat{\sigma}^2) \right |\right )\\
            &\leq \mathbb{E}\left ( \frac{N\epsilon\underline{\sigma}}{9N\underline{\sigma}}\sum_{i=1}^k\frac{1}{k\chi_{i,n(k,\des_i)-1}} \right )\\
            &\leq \frac{\epsilon}{3}.
        \end{split}
    \end{equation*}
    If event A occurs, 
    \begin{equation*}
        \begin{split}
             \frac{1}{k}\left | d^1(g_j,\hat{\mu},\hat{\sigma}^2)- d^1(\phi,\hat{\mu},\hat{\sigma}^2) \right |& \leq  \frac{N\epsilon\underline{\sigma}}{9N\underline{\sigma}}\sum_{i=1}^k\frac{1}{k\chi_{i,n(k,\des_i)-1}} \\
            & \leq \frac{\epsilon}{3}.
        \end{split}
    \end{equation*}
    Therefore, if event A occurs and $\left | d^1(g_j,\hat{\mu},\hat{\sigma}^2)-\mathbb{E}\left ( d^1(g_j,\hat{\mu},\hat{\sigma}^2)\right ) \right |/k\leq \epsilon/3$,
    \begin{equation*}
        \begin{split}
           \frac{1}{k}\left |d^1(\phi,\hat{\mu},\hat{\sigma}^2)-\mathbb{E}\left ( d^1(\phi,\hat{\mu},\hat{\sigma}^2)\right )\right |&\leq \frac{ \left |d^1(\phi,\hat{\mu},\hat{\sigma}^2)-d^1(g_j,\hat{\mu},\hat{\sigma}^2) \right |}{k} \\
           & +\frac{\left | d^1(g_j,\hat{\mu},\hat{\sigma}^2)-\mathbb{E}\left ( d^1(g_j,\hat{\mu},\hat{\sigma}^2)\right ) \right |}{k}\\
           & +\frac{\left |\mathbb{E}\left ( d^1(g_j,\hat{\mu},\hat{\sigma}^2)\right )-\mathbb{E}\left ( d^1(\phi,\hat{\mu},\hat{\sigma}^2)\right ) \right |}{k}\\
           &\leq \frac{\epsilon}{3}+\frac{\epsilon}{3}+\frac{\epsilon}{3}\\
           &= \epsilon.
        \end{split}
    \end{equation*}
    Thus, if event A occurs and $\max_{g_j}\left | d^1(g_j,\hat{\mu},\hat{\sigma}^2)-\mathbb{E}\left ( d^1(g_j,\hat{\mu},\hat{\sigma}^2)\right ) \right |/k\leq \epsilon/3$,
    $$\sup_{\phi \in \overline{M}^u_\lambda}\frac{1}{k}\left |d^1(\phi,\hat{\mu},\hat{\sigma}^2)-\mathbb{E}\left ( d^1(\phi,\hat{\mu},\hat{\sigma}^2)\right )\right | \leq \epsilon.$$
    Therefore, if $k\geq k^\star$, by the Union Bound,
    \begin{equation*}
        \begin{split}
            & \mathbb{P}\left (\sup_{\phi \in \overline{M}^u_\lambda}\frac{|d^1(\phi,\hat{\mu},\hat{\sigma}^2)-\mathbb{E}\left ( d^1(\phi,\hat{\mu},\hat{\sigma}^2)\right )|}{k}>\epsilon\right )\\
            &\leq \mathbb{P}\left ( A^C \bigcup \left \{\max_{g_j} \frac{\left | d^1(g_j,\hat{\mu},\hat{\sigma}^2)-\mathbb{E}\left ( d^1(g_j,\hat{\mu},\hat{\sigma}^2)\right ) \right |}{k}> \epsilon/3 \right \} \right )\\
            & \leq \mathbb{P}\left ( A^C\bigcup  \bigcup_{j=1}^t \left \{ \frac{\left | d^1(g_j,\hat{\mu},\hat{\sigma}^2)-\mathbb{E}\left ( d^1(g_j,\hat{\mu},\hat{\sigma}^2)\right ) \right |}{k}> \epsilon/3 \right \} \right )\\
            &\leq \mathbb{P}\left (A^C \right )+\sum_{j=1}^t \mathbb{P}\left ( \frac{\left | d^1(g_j,\hat{\mu},\hat{\sigma}^2)-\mathbb{E}\left ( d^1(g_j,\hat{\mu},\hat{\sigma}^2)\right ) \right |}{k}> \epsilon/3 \right )\\
            &=\frac{\epsilon}{2}+\sum_{j=1}^t \frac{\epsilon}{2t} \\
            &= \epsilon.
        \end{split}
    \end{equation*}
\end{proof}


\begin{thebibliography}{99}
\bibitem[durrett2019probability]{durrett2019probability} R.~Durrett, \emph{Probability: Theory and Examples}, 5th~ed.\hskip 1em plus 0.5em minus 0.4em\relax Cambridge University Press, 2019.
\bibitem[newey1991uniform]{newey1991uniform} W.~K. Newey, ``Uniform law of large numbers and empirical central limit theorems for time series models,'' \emph{Econometric Theory}, vol.~7, no.~4, pp. 431--467, 1991.
\end{thebibliography}
\end{document}
